% --- Archon physics formalization source begin ---
% archon:physics
% archon:covers PhyXMiniProblems/problem_phyx_mini_0127.lean
% archon:source-report reports/phyx_mini/problem_phyx_mini_0127.source.json

\chapter{Physics problem phyx\_mini\_0127}
\label{ch:PhyXMiniProblems_problem_phyx_mini_0127}

\paragraph{Problem source.}
\#\# Physical scenario

A soap bubble appears green \textbackslash{}(\textbackslash{}lambda=540\textbackslash{},\textbackslash{}mathrm\{nm\})\textbackslash{}) at the point on its front surface nearest the viewer. Assume \textbackslash{}(n=1.35\textbackslash{}).

\#\# Question

What is the smallest thickness the soap bubble film could have?

\#\# Answer choices

- A: A: 98nm

- B: B: 96nm

- C: C: 100nm

- D: D: 102nm

\#\# Figure caption (auxiliary; use the image as primary evidence)

The image is a diagram illustrating the interaction of light with a thin film, similar to a soap bubble. 

- **Objects and Features:**
  - A thin green film represents the bubble.
  - The film is labeled with a refractive index \textbackslash{}( n = 1.35 \textbackslash{}).
  - The film separates two regions: "Outside air" and "Bubble interior," each labeled with the refractive index \textbackslash{}( n = 1.00 \textbackslash{}).

- **Light Rays:**
  - An "Incident ray" approaches the bubble from the left.
  - "Reflected rays" diverge from the point of incidence on the film.

- **Labels and Text:**
  - "Incident ray" and "Reflected rays" are labeled with arrows to indicate direction.
  - The film thickness is indicated with \textbackslash{}( t \textbackslash{}) marked along the horizontal direction between dashed lines.

The diagram is demonstrating light reflection and refraction through a medium with a different refractive index than the surrounding air.

\#\# Dataset metadata

Wave Optics; Physical Model Grounding Reasoning, Spatial Relation Reasoning

\paragraph{Recorded answer/context.}
C: C: 100nm

\paragraph{Figure/image path.}
/root/proposal\_for\_physic/science-mango/phyx\_mini\_run/phyx\_data/test\_image/127.png

\paragraph{Formalization target.}
create a compiling Lean file with sorry bodies at `PhyXMiniProblems/problem\_phyx\_mini\_0127.lean`. The Lean declarations must preserve the physical quantities, dimensions or dimensional roles, figure labels, governing-law hypotheses, and final relation expressed by this problem.
Use Mathlib/Physlib names found through LeanExplore where available. If a physics API is missing, introduce faithful local abstractions rather than scalar placeholder aliases.

\begin{theorem}[Physics formalization target]
\label{thm:physics:phyx_mini_0127:target}
\lean{PhyXMiniProblems.ProblemPhyXMini0127.smallestGreenFilmThickness_isRecordedChoiceC}
\uses{def:physics:phyx-mini-0127:phyxminiproblems-problemphyxmini0127-nanometersvalue, def:physics:phyx-mini-0127:phyxminiproblems-problemphyxmini0127-soapbubblethinfilm, def:physics:phyx-mini-0127:phyxminiproblems-problemphyxmini0127-matchessoapbubblefigure, def:physics:phyx-mini-0127:phyxminiproblems-problemphyxmini0127-hasphysicalsoapfilmparameters, def:physics:phyx-mini-0127:phyxminiproblems-problemphyxmini0127-obeysthinfilmreflectionlaws, def:physics:phyx-mini-0127:phyxminiproblems-problemphyxmini0127-appearsgreenatnearestpoint, def:physics:phyx-mini-0127:phyxminiproblems-problemphyxmini0127-possiblegreenthicknessesnm, def:physics:phyx-mini-0127:phyxminiproblems-problemphyxmini0127-answerthicknessnanometers, def:physics:phyx-mini-0127:phyxminiproblems-problemphyxmini0127-recordedanswerchoice}
The assigned autoformalize agent should translate this physics problem into Lean declarations in the covered file, with theorem and lemma proof bodies written as `by sorry`.
\end{theorem}
\begin{proof}
This is an autoformalization task, not a proof task. Produce statements that can later be proved by the physics prover without weakening the physical model.
\end{proof}

% --- Archon declaration topology begin ---
\section{Declaration topology}
\begin{definition}[Dim Length]
  \label{def:physics:phyx-mini-0127:phyxminiproblems-problemphyxmini0127-dimlength}
  \lean{PhyXMiniProblems.ProblemPhyXMini0127.DimLength}
  A physical length with a real scalar carrier
\end{definition}
\begin{proof}
  This is the declared model definition of Dim Length.
\end{proof}
\begin{definition}[nanometer Unit Choices]
  \label{def:physics:phyx-mini-0127:phyxminiproblems-problemphyxmini0127-nanometerunitchoices}
  \lean{PhyXMiniProblems.ProblemPhyXMini0127.nanometerUnitChoices}
  Unit choices whose length unit is the nanometer used by the problem
\end{definition}
\begin{proof}
  This is the declared model definition of nanometer Unit Choices.
\end{proof}
\begin{definition}[nanometers Value]
  \label{def:physics:phyx-mini-0127:phyxminiproblems-problemphyxmini0127-nanometersvalue}
  \lean{PhyXMiniProblems.ProblemPhyXMini0127.nanometersValue}
  \uses{def:physics:phyx-mini-0127:phyxminiproblems-problemphyxmini0127-dimlength, def:physics:phyx-mini-0127:phyxminiproblems-problemphyxmini0127-nanometerunitchoices}
  The scalar nanometer readout of a physical length
\end{definition}
\begin{proof}
  This is the declared model definition of nanometers Value.
\end{proof}
\begin{definition}[Diagram Plane]
  \label{def:physics:phyx-mini-0127:phyxminiproblems-problemphyxmini0127-diagramplane}
  \lean{PhyXMiniProblems.ProblemPhyXMini0127.DiagramPlane}
  The Euclidean plane containing the thin-film diagram
\end{definition}
\begin{proof}
  This is the declared model definition of Diagram Plane.
\end{proof}
\begin{definition}[Optical Region]
  \label{def:physics:phyx-mini-0127:phyxminiproblems-problemphyxmini0127-opticalregion}
  \lean{PhyXMiniProblems.ProblemPhyXMini0127.OpticalRegion}
  The three physically distinct optical regions in the figure
\end{definition}
\begin{proof}
  This is the declared model definition of Optical Region.
\end{proof}
\begin{definition}[Film Interface]
  \label{def:physics:phyx-mini-0127:phyxminiproblems-problemphyxmini0127-filminterface}
  \lean{PhyXMiniProblems.ProblemPhyXMini0127.FilmInterface}
  The two film interfaces met by the incident light
\end{definition}
\begin{proof}
  This is the declared model definition of Film Interface.
\end{proof}
\begin{definition}[Reflected Ray Label]
  \label{def:physics:phyx-mini-0127:phyxminiproblems-problemphyxmini0127-reflectedraylabel}
  \lean{PhyXMiniProblems.ProblemPhyXMini0127.ReflectedRayLabel}
  Labels for the two reflected rays whose interference is observed
\end{definition}
\begin{proof}
  This is the declared model definition of Reflected Ray Label.
\end{proof}
\begin{definition}[incident Region]
  \label{def:physics:phyx-mini-0127:phyxminiproblems-problemphyxmini0127-incidentregion}
  \lean{PhyXMiniProblems.ProblemPhyXMini0127.incidentRegion}
  \uses{def:physics:phyx-mini-0127:phyxminiproblems-problemphyxmini0127-opticalregion, def:physics:phyx-mini-0127:phyxminiproblems-problemphyxmini0127-filminterface}
  The incident region of each oriented film interface
\end{definition}
\begin{proof}
  This is the declared model definition of incident Region.
\end{proof}
\begin{definition}[transmitted Region]
  \label{def:physics:phyx-mini-0127:phyxminiproblems-problemphyxmini0127-transmittedregion}
  \lean{PhyXMiniProblems.ProblemPhyXMini0127.transmittedRegion}
  \uses{def:physics:phyx-mini-0127:phyxminiproblems-problemphyxmini0127-opticalregion, def:physics:phyx-mini-0127:phyxminiproblems-problemphyxmini0127-filminterface}
  The transmitted region of each oriented film interface
\end{definition}
\begin{proof}
  This is the declared model definition of transmitted Region.
\end{proof}
\begin{definition}[Ray Segment]
  \label{def:physics:phyx-mini-0127:phyxminiproblems-problemphyxmini0127-raysegment}
  \lean{PhyXMiniProblems.ProblemPhyXMini0127.RaySegment}
  \uses{def:physics:phyx-mini-0127:phyxminiproblems-problemphyxmini0127-diagramplane}
  A directed straight segment used to retain the ray labels in the figure
\end{definition}
\begin{proof}
  This is the declared model definition of Ray Segment.
\end{proof}
\begin{definition}[Inner Surface Reflected Route]
  \label{def:physics:phyx-mini-0127:phyxminiproblems-problemphyxmini0127-innersurfacereflectedroute}
  \lean{PhyXMiniProblems.ProblemPhyXMini0127.InnerSurfaceReflectedRoute}
  \uses{def:physics:phyx-mini-0127:phyxminiproblems-problemphyxmini0127-diagramplane}
  The ray reflected at the inner film surface: it enters the film at the front surface, reflects at the film--interior interface, returns to the front surface, and then travels back toward the viewer
\end{definition}
\begin{proof}
  This is the declared model definition of Inner Surface Reflected Route.
\end{proof}
\begin{definition}[Soap Bubble Thin Film]
  \label{def:physics:phyx-mini-0127:phyxminiproblems-problemphyxmini0127-soapbubblethinfilm}
  \lean{PhyXMiniProblems.ProblemPhyXMini0127.SoapBubbleThinFilm}
  \uses{def:physics:phyx-mini-0127:phyxminiproblems-problemphyxmini0127-dimlength, def:physics:phyx-mini-0127:phyxminiproblems-problemphyxmini0127-diagramplane, def:physics:phyx-mini-0127:phyxminiproblems-problemphyxmini0127-opticalregion, def:physics:phyx-mini-0127:phyxminiproblems-problemphyxmini0127-reflectedraylabel, def:physics:phyx-mini-0127:phyxminiproblems-problemphyxmini0127-raysegment, def:physics:phyx-mini-0127:phyxminiproblems-problemphyxmini0127-innersurfacereflectedroute}
  The physical quantities and labeled geometry of the depicted soap bubble. constructivelyReflectsGreenAtThicknessNm x is the observable reflected-light condition for a hypothetical film with nanometer thickness readout x. Its dependence on optical physics is supplied separately by ObeysThinFilmReflectionLaws; it is not defined to make the target true
\end{definition}
\begin{proof}
  This is the declared model definition of Soap Bubble Thin Film.
\end{proof}
\begin{definition}[Matches Soap Bubble Figure]
  \label{def:physics:phyx-mini-0127:phyxminiproblems-problemphyxmini0127-matchessoapbubblefigure}
  \lean{PhyXMiniProblems.ProblemPhyXMini0127.MatchesSoapBubbleFigure}
  \uses{def:physics:phyx-mini-0127:phyxminiproblems-problemphyxmini0127-nanometersvalue, def:physics:phyx-mini-0127:phyxminiproblems-problemphyxmini0127-soapbubblethinfilm}
  Figure and problem readouts: green light has wavelength 540 nm, the soap film has index 1.35, and both adjacent regions are air with index 1.00. The viewer-nearest surface point gives normal incidence in this local model. The route equalities retain the incident and two reflected rays shown in the source image. No numerical value of the unknown thickness occurs here
\end{definition}
\begin{proof}
  This is the declared model definition of Matches Soap Bubble Figure.
\end{proof}
\begin{definition}[Has Physical Soap Film Parameters]
  \label{def:physics:phyx-mini-0127:phyxminiproblems-problemphyxmini0127-hasphysicalsoapfilmparameters}
  \lean{PhyXMiniProblems.ProblemPhyXMini0127.HasPhysicalSoapFilmParameters}
  \uses{def:physics:phyx-mini-0127:phyxminiproblems-problemphyxmini0127-nanometersvalue, def:physics:phyx-mini-0127:phyxminiproblems-problemphyxmini0127-incidentregion, def:physics:phyx-mini-0127:phyxminiproblems-problemphyxmini0127-transmittedregion, def:physics:phyx-mini-0127:phyxminiproblems-problemphyxmini0127-soapbubblethinfilm}
  Positivity and optical-density ordering for an ordinary soap film surrounded by air. These conditions contain no value for the requested minimum
\end{definition}
\begin{proof}
  This is the declared model definition of Has Physical Soap Film Parameters.
\end{proof}
\begin{definition}[Obeys Thin Film Reflection Laws]
  \label{def:physics:phyx-mini-0127:phyxminiproblems-problemphyxmini0127-obeysthinfilmreflectionlaws}
  \lean{PhyXMiniProblems.ProblemPhyXMini0127.ObeysThinFilmReflectionLaws}
  \uses{def:physics:phyx-mini-0127:phyxminiproblems-problemphyxmini0127-nanometersvalue, def:physics:phyx-mini-0127:phyxminiproblems-problemphyxmini0127-soapbubblethinfilm}
  The governing normal-incidence reflected-light laws for an air--film--air layer. Reflection from outside air into the denser film contributes a half-turn phase reversal, while reflection from film to interior air does not. Consequently constructive reflection occurs exactly when the round-trip optical thickness is an odd half-wavelength: 4 n\_film t = (2 m + 1) λ for some nonnegative interference order m.
\end{definition}
\begin{proof}
  This is the declared model definition of Obeys Thin Film Reflection Laws.
\end{proof}
\begin{definition}[Appears Green At Nearest Point]
  \label{def:physics:phyx-mini-0127:phyxminiproblems-problemphyxmini0127-appearsgreenatnearestpoint}
  \lean{PhyXMiniProblems.ProblemPhyXMini0127.AppearsGreenAtNearestPoint}
  \uses{def:physics:phyx-mini-0127:phyxminiproblems-problemphyxmini0127-nanometersvalue, def:physics:phyx-mini-0127:phyxminiproblems-problemphyxmini0127-soapbubblethinfilm}
  The observed bubble thickness produces constructive green reflection
\end{definition}
\begin{proof}
  This is the declared model definition of Appears Green At Nearest Point.
\end{proof}
\begin{definition}[possible Green Thicknesses Nm]
  \label{def:physics:phyx-mini-0127:phyxminiproblems-problemphyxmini0127-possiblegreenthicknessesnm}
  \lean{PhyXMiniProblems.ProblemPhyXMini0127.possibleGreenThicknessesNm}
  \uses{def:physics:phyx-mini-0127:phyxminiproblems-problemphyxmini0127-soapbubblethinfilm}
  Positive nanometer thicknesses that could produce the observed green color
\end{definition}
\begin{proof}
  This is the declared model definition of possible Green Thicknesses Nm.
\end{proof}
\begin{definition}[Answer Choice]
  \label{def:physics:phyx-mini-0127:phyxminiproblems-problemphyxmini0127-answerchoice}
  \lean{PhyXMiniProblems.ProblemPhyXMini0127.AnswerChoice}
  Labels of the four displayed multiple-choice answers
\end{definition}
\begin{proof}
  This is the declared model definition of Answer Choice.
\end{proof}
\begin{definition}[answer Thickness Nanometers]
  \label{def:physics:phyx-mini-0127:phyxminiproblems-problemphyxmini0127-answerthicknessnanometers}
  \lean{PhyXMiniProblems.ProblemPhyXMini0127.answerThicknessNanometers}
  \uses{def:physics:phyx-mini-0127:phyxminiproblems-problemphyxmini0127-answerchoice}
  Film-thickness readout in nanometers printed beside each answer choice
\end{definition}
\begin{proof}
  This is the declared model definition of answer Thickness Nanometers.
\end{proof}
\begin{definition}[recorded Answer Choice]
  \label{def:physics:phyx-mini-0127:phyxminiproblems-problemphyxmini0127-recordedanswerchoice}
  \lean{PhyXMiniProblems.ProblemPhyXMini0127.recordedAnswerChoice}
  \uses{def:physics:phyx-mini-0127:phyxminiproblems-problemphyxmini0127-answerchoice}
  The answer label recorded by the source dataset
\end{definition}
\begin{proof}
  This is the declared model definition of recorded Answer Choice.
\end{proof}
% --- Archon declaration topology end ---
% --- Archon physics formalization source end ---
